\subsection{COMMUTATION LEMMAS}

Lemma 1. Let \(\mathbf{A}\) be a commutative ring, \(\mathbf{E}\) an \(\mathbf{A}\) -algebra, \((x_{i})_{1\leqslant i\leqslant n}\) a finite sequence of elements of \(\mathbf{E}\) , \((\lambda_i)_{1\leqslant i\leqslant n}\) a finite sequence of elements of \(\mathbf{A}\) and \(y\) an element of \(\mathbf{E}\) ; suppose that

\[
x _ {i} y = \lambda_ {i} y x _ {i} \quad \text { for } 1 \leqslant i \leqslant n.\tag{13}
\]

Then

\[
\left(x _ {1} x _ {2} \dots x _ {n}\right) y = \left(\lambda_ {1} \lambda_ {2} \dots \lambda_ {n}\right) y \left(x _ {1} x _ {2} \dots x _ {n}\right).\tag{14}
\]

The lemma being trivial for \(n = 1\) , we argue by induction on \(n \geqslant 2\) . Now

\[
\begin{array}{r l} (x _ {1} x _ {2} \dots x _ {n}) y & = (x _ {1} \dots x _ {n - 1}) (x _ {n} y) \\ & = (x _ {1} \dots x _ {n - 1}) (\lambda_ {n} y x _ {n}) = \lambda_ {n} ((x _ {1} \dots x _ {n - 1}) y) x _ {n}, \end{array}
\]

which, by the induction hypothesis, is equal to

\[
\lambda_ {n} \left(\lambda_ {1} \dots \lambda_ {n - 1}\right) y \left(x _ {1} \dots x _ {n - 1}\right) x _ {n} = \left(\lambda_ {1} \dots \lambda_ {n - 1} \lambda_ {n}\right) y \left(x _ {1} \dots x _ {n - 1} x _ {n}\right),
\]

whence the lemma.

Lemma 2. Let A be a commutative ring, E an A-algebra and \((x_{i})_{1\leqslant i\leqslant n}\) and \((y_{i})_{1\leqslant i\leqslant n}\) two finite sequences of \(n\) elements of E; suppose that for \(1\leqslant j\leqslant i\leqslant n\) ,

\[
x _ {i} y _ {j} = \lambda_ {i j} y _ {j} x _ {i} \quad \text { with } \lambda_ {i j} \in A.\tag{15}
\]

Then

\[
\left(x _ {1} x _ {2} \dots x _ {n}\right) \left(y _ {1} y _ {2} \dots y _ {n}\right) = \left(\prod_ {i > j} \lambda_ {i j}\right) \left(x _ {1} y _ {1}\right) \left(x _ {2} y _ {2}\right) \dots \left(x _ {n} y _ {n}\right).\tag{16}
\]

The lemma being trivial for \(n = 1\) , we again argue by induction on \(n\) for \(n \geqslant 2\) . By virtue of Lemma 1,

\[
\begin{array}{r l} (x _ {1} \dots x _ {n}) (y _ {1} \dots y _ {n}) & = x _ {1} (x _ {2} \dots x _ {n}) y _ {1} (y _ {2} \dots y _ {n}) \\ & = \left(\prod_ {i > 1} \lambda_ {i 1}\right) (x _ {1} y _ {1}) (x _ {2} \dots x _ {n}) (y _ {2} \dots y _ {n}) \end{array}
\]

and it then suffices to apply the induction hypothesis to obtain (16).

For every family \(\lambda = (\lambda_{ij})\) of elements of A, with \(1 \leqslant j < i \leqslant n\) , and for every permutation \(\sigma \in \mathfrak{S}_n\) , we write

\[
\varepsilon_ {\sigma} (\lambda) = \prod_ {i > j, \sigma^ {- 1} (i) < \sigma^ {- 1} (j)} \lambda_ {i j} = \prod_ {i <   j, \sigma (i) > \sigma (j)} \lambda_ {\sigma (i), \sigma (j)}.\tag{17}
\]

Observe that, when \(\mathbf{A} = \mathbf{Z}\) and \(\lambda_{ij} = -1\) for every ordered pair \((i,j)\) such that \(1 \leqslant j < i \leqslant n\) , \(\varepsilon_{\sigma}(\lambda)\) is just the signature \(\varepsilon_{\sigma}\) of the permutation \(\sigma\) (I, § 5, no. 7).

Lemma 3. Let A be a commutative ring, E an A-algebra, \((x_{i})_{1\leqslant i\leqslant n}\) a finite sequence of elements of E and suppose that, for every ordered pair \((i,j)\) of integers such that \(1\leqslant j < i\leqslant n\) ,

\[
x _ {i} x _ {j} = \lambda_ {i j} x _ {j} x _ {i} \quad \text { with } \lambda_ {i j} \in A.\tag{18}
\]

Then, for every permutation \(\sigma \in \mathfrak{S}_n\) ,

\[
x _ {\sigma (1)} x _ {\sigma (2)} \dots x _ {\sigma (n)} = \varepsilon_ {\sigma} (\lambda) x _ {1} x _ {2} \dots x _ {n}.\tag{19}
\]

The lemma is trivial for \(n = 1\) and \(n = 2\) ; we proceed by induction on \(n\) for \(n \geqslant 3\) . If \(\sigma(n) = n\) , relation (19) follows from the induction hypothesis. Suppose therefore that \(\sigma(n) = k, k \neq n\) , and let \(\tau\) be the permutation of \([1, n]\) defined by

\[
\left\{ \begin{array}{l l} \tau (i) = i & \text { for } i <   k \\ \tau (i) = i + 1 & \text { for } k \leqslant i <   n \\ \tau (n) = k. \end{array} \right.
\]

Let \(\pi = \tau^{-1} \circ \sigma\) ; the permutation \(\pi\) leaves \(n\) fixed; now \(\sigma = \tau \circ \pi\) and therefore, writing \(y_i = x_{\tau(i)}\) , \(y_{\pi(i)} = x_{\sigma(i)}\) . If \(i \neq n\) and \(j \neq n\) , the relations \(\pi(i) > \pi(j)\) and \(\sigma(i) > \sigma(j)\) are equivalent (since \(\tau\) is a strictly increasing mapping of \([1, n - 1]\) into \([1, n]\) . For \(i \neq n\) , \(j \neq n\) and \(\sigma(i) > \sigma(j)\) ,

\[
y _ {\pi (i)} y _ {\pi (j)} = x _ {\sigma (i)} x _ {\sigma (j)} = \lambda_ {\sigma (i), \sigma (j)} x _ {\sigma (j)} x _ {\sigma (i)} = \lambda_ {\sigma (i), \sigma (j)} y _ {\pi (j)} y _ {\pi (i)}
\]

whence, by the induction hypothesis (using the fact that \(\pi(n) = n\) ):

\[
y _ {\pi (1)} y _ {\pi (2)} \dots y _ {\pi (n)} = \left(\prod_ {i <   j <   n, \sigma (i) > \sigma (j)} \lambda_ {\sigma (i), \sigma (j)}\right) y _ {1} y _ {2} \dots y _ {n}
\]

that is

\[
x _ {\sigma (1)} x _ {\sigma (2)} \dots x _ {\sigma (n)} = \left(\prod_ {i <   j <   n, \sigma (i) > \sigma (j)} \lambda_ {\sigma (i), \sigma (j)}\right) x _ {\tau (1)} \dots x _ {\tau (n)}.\tag{20}
\]

Now

\[
x _ {\tau (1)} \dots x _ {\tau (n)} = x _ {1} \dots x _ {k - 1} x _ {k + 1} \dots x _ {n} x _ {k}
\]

and this, by Lemma 1, is equal to

\[
\left(\prod_ {j > k} \lambda_ {j k}\right) x _ {1} \dots x _ {n} = \left(\prod_ {\sigma (i) > \sigma (n)} \lambda_ {\sigma (i), \sigma (n)}\right) x _ {1} \dots x _ {n}.\tag{21}
\]

Finally, (20) and (21) give

\[
x _ {\sigma (1)} \dots x _ {\sigma (n)} = \alpha . x _ {1} \dots x _ {n}
\]

with

\[
\begin{array}{l} \alpha = \left(\prod_ {i <   j <   n, \sigma (i) > \sigma (j)} \lambda_ {\sigma (i), \sigma (j)}\right) \cdot \left(\prod_ {i <   n, \sigma (i) > \sigma (n)} \lambda_ {\sigma (i), \sigma (n)}\right) \\ = \prod_ {i <   j, \sigma (i) > \sigma (j)} \lambda_ {\sigma (i), \sigma (j)} = \varepsilon_ {\sigma} (\lambda) \end{array}
\]

which completes the proof of Lemma 3.

